\newpage
\section{Tiempo hasta velocidad vertical cero con resistencia del aire}
Durante el ascenso, la computadora utilizará la velocidad estimada $V_{r,k}$ y la aceleración vertical de trayectoria obtenida del acelerómetro. No se impondrá $a_z=-g$ como valor permanente: la medición permite incorporar el efecto actual del empuje y de las fuerzas aerodinámicas. La estimación requiere lecturas válidas, calibradas y corregidas por orientación.

\subsection{Qué mide el acelerómetro durante el vuelo}
Un acelerómetro mide fuerza específica, no la aceleración gravitatoria directamente. En caída libre ideal, sin resistencia del aire, su lectura es aproximadamente cero aunque el cohete acelere hacia abajo a $g$. En reposo sobre una superficie, la fuerza de soporte produce una lectura de magnitud cercana a $g$.\footnote{Analog Devices, ficha técnica ADXL366, explicación de detección de caída libre: \href{https://www.analog.com/media/en/technical-documentation/data-sheets/adxl366.pdf}{documentación del fabricante}. Se cita el principio físico, no el modelo de sensor del proyecto.}

Definimos la componente vertical local de fuerza específica como
\begin{equation}
 f_{z,N,k}=\mathbf e_z^{\mathsf T}\mathbf R_{NB,k}
 (\mathbf f_{m,k}-\mathbf b_f),
 \qquad \boxed{a_{z,k}=f_{z,N,k}-g}.
\end{equation}
La lectura se supone corregida en escala y expresada en $\mathrm{m/s^2}$. Si el cohete está inclinado, $f_{z,N}$ combina sus tres ejes: no debe sustituirse por la componente $z$ cruda del sensor. Para la rotación lateral del ejemplo,
\begin{equation}
 a_z=-f_x\sin\theta+f_z\cos\theta-g.
\end{equation}
En un modelo de traslación referido al centro de masa,
\begin{equation}
 f_{z,N}=\frac{F_{T,z}+F_{\mathrm{aero},z}}{m},\qquad
 a_z=\frac{F_{T,z}+F_{\mathrm{aero},z}}{m}-g.
\end{equation}
$F_{T,z}$ es el empuje vertical y $F_{\mathrm{aero},z}$ la resultante aerodinámica vertical. Después del empuje, el sensor sigue registrando el efecto de la aerodinámica. Durante ascenso en aire quieto y con solo arrastre, ese efecto es negativo y $a_z$ es más negativa que $-g$.

\subsection{Cómo expresar el efecto del viento y el arrastre}
El arrastre depende de la velocidad respecto al aire. Si $\mathbf v_N$ es la velocidad del cohete respecto al suelo y $\mathbf w_N$ la velocidad del viento,
\begin{equation}
 \mathbf u=\mathbf v_N-\mathbf w_N,\qquad
 \mathbf F_D=-\tfrac12\rho C_D A\,\|\mathbf u\|\mathbf u.
\end{equation}
Aquí $\rho$ es la densidad del aire, $C_D$ el coeficiente de arrastre y $A$ su área de referencia. Es un modelo de arrastre opuesto al flujo relativo; una fuerza de sustentación o normal al flujo requeriría un término adicional.\footnote{NASA Glenn Research Center: \href{https://www1.grc.nasa.gov/beginners-guide-to-aeronautics/modern-drag-equation/}{Modern Drag Equation}.}
La componente vertical del modelo es
\begin{equation}
 a_z=\frac{F_{T,z}}{m}-g-
 \frac{\rho C_D A}{2m}\|\mathbf v_N-\mathbf w_N\|(V_z-w_z).
\end{equation}
El viento lateral modifica la rapidez relativa y, con ello, el arrastre; el viento vertical también puede cambiar el signo de su componente vertical. La sola estimación $V_r$ no determina el viento ni toda la velocidad relativa.

\newpage
\subsection{Predicción con la aceleración vertical medida}
Para actualizar $V_r$ se utilizará la aceleración vertical corregida durante el vuelo. Para predecir el apogeo, además se necesita estar ascendiendo y desacelerando. Sea $\bar a_{z,k}$ una estimación filtrada de esa aceleración, calculada con mediciones recientes; la barra indica suavizado temporal, no conocimiento de la aceleración futura.

Como aproximación local de movimiento rectilíneo uniformemente acelerado (MRUA), se supone que $\bar a_{z,k}$ permanece constante durante la predicción:
\begin{equation}
 V_z(t_k+\Delta t)\approx V_{r,k}+\bar a_{z,k}\Delta t.
\end{equation}
Al imponer velocidad vertical cero,
\begin{equation}
 \boxed{\Delta t_0\approx-\frac{V_{r,k}}{\bar a_{z,k}}},
 \qquad V_{r,k}>0,\quad \bar a_{z,k}<0.
\end{equation}
$\Delta t_0$ es el tiempo restante, mientras que $t_k+\Delta t_0$ es el instante estimado de apogeo. Si la aceleración es positiva, el modelo no predice una detención futura durante el ascenso. Si es cero o muy cercana a cero, no proporciona un tiempo finito fiable. En las lecturas del ejemplo anterior, $a_{z,2}=+2.662\,\mathrm{m/s^2}$, por lo que todavía no corresponde aplicar esta predicción de apogeo.

\subsection{Ejemplo después de terminar el empuje}
Supongamos ahora un instante distinto, de ascenso sin empuje, con $V_r=8.328\,\mathrm{m/s}$ y una fuerza específica vertical local filtrada de $-2.00\,\mathrm{m/s^2}$. Atribuyendo esta última al arrastre,
\begin{equation}
 \bar a_z=-2.00-9.81=-11.81\,\mathrm{m/s^2}.
\end{equation}
La estimación local del tiempo restante es
\begin{equation}
 \Delta t_0\approx\frac{8.328}{11.81}\approx0.705\,\mathrm{s}.
\end{equation}
Si se ignorara el arrastre se obtendría $8.328/9.81\approx0.849\,\mathrm{s}$. La diferencia refleja la desaceleración adicional que sí aparece en la medición. No debemos volver a restar un arrastre calculado a $\bar a_z$: hacerlo contaría dos veces la misma fuerza. El modelo aerodinámico sirve para interpretar o predecir la evolución, no para duplicar una contribución ya medida.

\subsection{Alcance de la estimación}
El acelerómetro observa el efecto actual del aire, pero no determina cómo cambiará hasta el apogeo. En ascenso puramente vertical, sin viento, sin empuje y con masa y parámetros aerodinámicos constantes, el modelo es
\begin{equation}
 \dot V_z=-g-\beta V_z^2,\qquad \beta=\frac{\rho C_D A}{2m}.
\end{equation}
Al disminuir la velocidad, disminuye el arrastre: mantener la desaceleración inicial constante tiende a subestimar el tiempo restante en este caso. Por eso $0.705\,\mathrm{s}$ es una extrapolación local, no un tiempo exacto. Se recalculará con cada actualización válida, considerando el retardo del filtrado.

Si el sensor está alejado del centro de masa, su lectura incluye aceleraciones tangenciales y centrípetas por la rotación; deben compensarse o demostrarse despreciables. Las lecturas saturadas o la orientación incierta invalidan la predicción. Finalmente, velocidad vertical cero no implica velocidad lateral cero, y la predicción temporal no sustituye la confirmación medida del apogeo.
\end{document}
